% ===================================================================
%  QUADRO N.º 16 — Os Grandes Armazéns. — O Bazar. — Os Brinquedos.
%
%  Ceci n'est PAS une transcription : c'est une traduction, faite en
%  2026, en portugais d'aujourd'hui. Voir texto/pt/10-quadro-01.tex
%  pour la consigne, qui vaut pour les seize fichiers.
%
%  Au 15, l'ido saute la lettre (d) du tableau d'astronomie ; la
%  traduction la garde, comme les sept autres colonnes, et
%  outils/renvoji.py la declare en ecart connu.
% ===================================================================

%%K t16-apar-1 apar
\VUsaut{4.0mm}
\VUtitre{15.0pt}{60}{QUADRO N.º 16}
\VUsaut{2.0mm}
\VUpk{13.2pt}{40}{Os Grandes Armazéns. --- O Bazar.}
\VUpk{13.2pt}{40}{Os Brinquedos.}
\VUsaut{3.4mm}

%%K t16-01-1 p
1. --- Aproxima-se o ano novo. Os grandes armazéns prepararam as suas
exposições de \VUgras{brinquedos} \textsuperscript{(1)} e presentes de ano novo.

%%K t16-01-2 p
Pedi ao meu avô que me levasse ao bazar a ver essa bela exposição, e ele
aceitou.

%%K t16-02-1 p
2. --- Primeiro parámos diante de umas belas panóplias de armas com uma
\VUgras{espingarda}, um \VUgras{boné de plaqueta}, um \VUgras{sabre}, um
\VUgras{boldrié} e um par de \VUgras{dragonas}, ou então um \VUgras{capacete},
uma \VUgras{couraça} \textsuperscript{(3)} e uma \VUgras{pistola} \textsuperscript{(4)}.
Tudo isso me interessava muito mais do que as \VUgras{lanternas mágicas}
\textsuperscript{(2)}.

%%K t16-tit-1 sub
\VUsaut{3.0mm}
\VUpk{10.2pt}{40}{Belos espetáculos.}
\VUsaut{2.6mm}

%%K t16-03-1 p
3. --- Na mesma secção havia uma \VUgras{sentinela} \textsuperscript{(5)} na sua
\VUgras{guarita}; parecia montar guarda sobre toda uma casa de feras de
cartão. Quantos animais havia! Um \VUgras{papagaio} \textsuperscript{(6)} no seu
poleiro, um \VUgras{rinoceronte} \textsuperscript{(7)} com um corno no nariz, uma
\VUgras{rena} \textsuperscript{(8)}, uma \VUgras{avestruz} \textsuperscript{(9)}, um
\VUgras{macaco} \textsuperscript{(10)} que partia uma noz, uma \VUgras{raposa}
\textsuperscript{(11)} que levava uma galinha, um \VUgras{crocodilo}
\textsuperscript{(12)} com as suas enormes fauces abertas, um \VUgras{camelo}
\textsuperscript{(13)}, um elegante \VUgras{galgo} \textsuperscript{(14)}, uma
\VUgras{pantera} \textsuperscript{(15)}, e depois dois bichos que eu não conhecia,
que são, segundo me dizem, uma \VUgras{doninha} \textsuperscript{(16)} e uma
\VUgras{hiena} \textsuperscript{(17)}. Havia também um \VUgras{leopardo}
\textsuperscript{(18)} e um \VUgras{lobo} \textsuperscript{(21)}; muito perto havia
lindas \VUgras{ratazanas} \textsuperscript{(19)} e \VUgras{ratos}
\textsuperscript{(20)} minúsculos de borracha. Por fim, um \VUgras{hipopótamo}
\textsuperscript{(22)} e um \VUgras{dromedário} \textsuperscript{(23)} corcovado
completavam a coleção. Ter-me-ia ficado ali horas inteiras se não
houvesse, muito perto, um esplêndido acampamento cheio de
\VUgras{soldadinhos de chumbo} \textsuperscript{(29)} que me chamou a atenção.

%%K t16-tit-2 sub
\VUsaut{4.0mm}
\VUpk{10.2pt}{40}{Soldados e guerra.}
\VUsaut{2.8mm}

%%K t16-04-1 p
4. --- O meu avô, que é um \VUgras{inválido de guerra} \textsuperscript{(24)} e
teve de deixar o exército por causa de um \VUgras{ferimento} que lhe valeu
a \VUgras{medalha militar}, adora contar-me as \VUgras{campanhas} em que
participou. Encantou-o explicar-me os diferentes movimentos do pequeno
exército em miniatura. Um grupo de soldados a sério que visitavam o bazar
--- um \VUgras{engenheiro militar} \textsuperscript{(25)}, um \VUgras{caçador
alpino} \textsuperscript{(26)} com a sua \VUgras{boina}, um \VUgras{primeiro-sargento}
\textsuperscript{(27)} de infantaria e um \VUgras{sargento} \textsuperscript{(28)} de
artilharia com um \VUgras{galão} de ouro na manga --- deteve-se perto de
nós a escutar as explicações.

%%K t16-05-1 p
5. --- O meu avô, muito orgulhoso de ter um público tão distinto,
improvisou todo um plano de batalha.

%%K t16-05-2 p
A praça-forte, com as suas \VUgras{muralhas} e os seus \VUgras{fossos},
estava sitiada pelo \VUgras{exército inimigo}, que, após um longo
\VUgras{bombardeamento}, se dispunha a tomá-la de \VUgras{assalto}. Mas os
\VUgras{sitiados}, empurrados pela fome, tinham tentado uma \VUgras{saída}
contra o \VUgras{acampamento} dos \VUgras{sitiantes}. Repelido o
\VUgras{ataque} num encarniçado \VUgras{combate} em torno das
\VUgras{tendas}, os sitiantes \VUgras{vencedores} lançaram os seus
\VUgras{esquadrões} em \VUgras{perseguição} dos assaltantes
\VUgras{vencidos}. Estes \VUgras{retiraram-se}, deixando o \VUgras{campo de
batalha} semeado de \VUgras{mortos} e \VUgras{feridos}. Os
\VUgras{maqueiros} levaram os feridos ao \VUgras{hospital de campanha} para
que o \VUgras{cirurgião} os tratasse.

%%K t16-05-3 p
Apesar da heroica resistência de alguns \VUgras{batalhões} que se tinham
formado em \VUgras{quadrado}, a \VUgras{derrota} foi completa; o resto do
exército \VUgras{fugiu}, deixando atrás muitos \VUgras{prisioneiros}. O
inimigo estava prestes a coroar a sua \VUgras{vitória} tomando a praça.
Talvez fosse aquele o fim da campanha, e após um \VUgras{armistício}
poderia assinar-se um \VUgras{tratado de paz}.

%%K t16-tit-3 sub
\VUsaut{4.4mm}
\VUpk{10.2pt}{40}{Os presentes de ano novo.}
\VUsaut{2.8mm}

%%K t16-06-1 p
6. --- Os jovens soldados, que se riam a escutar o pitoresco relato do
veterano, fizeram-lhe mais algumas perguntas. Quanto a mim, dei a volta ao
balcão para ver uma bela \VUgras{besta} \textsuperscript{(32)} e um \VUgras{arco}
\textsuperscript{(33)} com a sua \VUgras{flecha} \textsuperscript{(31)}. Ali encontrei
outra coleção de animais: dois \VUgras{pássaros cantores} \textsuperscript{(34)}, um
\VUgras{rouxinol} e uma \VUgras{toutinegra} de corda que cantavam a plenos
pulmões, e uma série de bichos estranhos --- um \VUgras{morcego}
\textsuperscript{(35)}, uma \VUgras{jiboia} \textsuperscript{(36)}, uma \VUgras{tartaruga}
\textsuperscript{(37)}, uma \VUgras{foca} \textsuperscript{(38)}, uma \VUgras{baleia}
\textsuperscript{(39)} e um \VUgras{tubarão} \textsuperscript{(40)} feitos de pele de
batedor de ouro.

%%K t16-07-1 p
7. --- Não me detive muito a olhá-los, pois tinha visto o que sonhava: um
magnífico \VUgras{teatro de fantoches} \textsuperscript{(41)} com esplêndidos
\VUgras{cenários} e um delicioso \VUgras{pano} vermelho. No \VUgras{palco},
dois atores pareciam representar um \VUgras{papel} numa \VUgras{peça}
interessantíssima, pois os pequenos \VUgras{espetadores} de cartão dos
camarotes, ou sentados nas \VUgras{plateias} e na \VUgras{galeria} atrás do
\VUgras{maestro}, aplaudiam com entusiasmo.

%%K t16-07-2 p
Infelizmente, aquele teatro era caríssimo.

%%K t16-08-1 p
8. --- Além disso, era difícil escolher entre os brinquedos, pois as
montras estavam cheias de brinquedos de todas as espécies:
\VUgras{Arlequim} \textsuperscript{(42)}, \VUgras{Polichinelo} \textsuperscript{(43)},
\VUgras{Pierrô} \textsuperscript{(44)}, \VUgras{corujas} \textsuperscript{(45)} que batiam
as asas, \VUgras{melros} \textsuperscript{(46)} que assobiavam, \VUgras{cães de
água} que ladravam, um \VUgras{gramofone} \textsuperscript{(47)} e
\VUgras{quebra-cabeças}. Um desses brinquedos divertiu-me muitíssimo:
representava um \VUgras{combate naval} \textsuperscript{(48)} no qual um
\VUgras{navio} é atacado por \VUgras{piratas} em frente de uma costa
plantada de \VUgras{coqueiros}.

%%K t16-noto-1 noto
\VUnotes{143.2mm}{%
\textit{(*) Ver o quadro n.º 6.}%
}

%%K t16-09-1 p
9. --- Pude ver todas estas maravilhas com toda a calma, pois havia pouca
gente nos armazéns --- apenas um punhado de curiosos, um \VUgras{dragão}
\textsuperscript{(49)} e um \VUgras{cobrador} \textsuperscript{(50)} que apresentava uma
\VUgras{letra de câmbio} na \VUgras{caixa} \textsuperscript{(51)} principal.

%%K t16-10-1 p
10. --- Perguntei ao meu avô para que serviam os livros que havia nas
prateleiras atrás do \VUgras{caixa}; disse-me que eram os \VUgras{livros de
contas}.

%%K t16-tit-4 sub
\VUsaut{3.6mm}
\VUpk{10.2pt}{40}{A ver por toda a loja.}
\VUsaut{2.8mm}

%%K t16-11-1 p
11. --- Ao sair da secção dos brinquedos, fomos à selaria dar uma vista de
olhos às \VUgras{selas} \textsuperscript{(53)}, aos \VUgras{estribos}
\textsuperscript{(54)}, às \VUgras{bridas} \textsuperscript{(55)}, aos \VUgras{freios}
\textsuperscript{(56)} e aos \VUgras{chicotes}. Enquanto o \VUgras{chefe de secção}
\textsuperscript{(57)} dava ao meu avô algumas informações, fui ver a exposição
de \VUgras{porcelana} e \VUgras{cristais} \textsuperscript{(58)}, onde havia belas
\VUgras{saladeiras} e delicados \VUgras{bules} \textsuperscript{(59)}.

%%K t16-12-1 p
12. --- Para subir ao primeiro andar, passámos por trás da montra dos
\VUgras{espartilhos} \textsuperscript{(60)}, junto à retrosaria, onde estavam
expostos belíssimos \VUgras{calções} \textsuperscript{(63)}. Perto, uma
\VUgras{caixeira} \textsuperscript{(61)} ajudava uma \VUgras{cliente}
\textsuperscript{(62)} a escolher belos \VUgras{tecidos} \textsuperscript{(64)} de seda.

%%K t16-12-2 p
Ao subir a escada, reparei que os armazéns estavam enfeitados com
\VUgras{lanternas} \textsuperscript{(65)} japonesas, \VUgras{sombrinhas}
\textsuperscript{(66)} e enormes \VUgras{palmeiras} \textsuperscript{(70)}.

%%K t16-13-1 p
13. --- No primeiro andar não havia muito que ver: só móveis (*),
\VUgras{cofres-fortes} \textsuperscript{(67)}, \VUgras{cómodas} \textsuperscript{(68)} e
\VUgras{carrinhos de bebé} \textsuperscript{(69)}. Mas a vista geral dos armazéns
era muito \VUgras{entretida}.

%%K t16-14-1 p
14. --- Por baixo de nós via os instrumentos de música: \VUgras{harpas}
\textsuperscript{(71)}, \VUgras{guitarras} \textsuperscript{(72)}, \VUgras{bandolins}
\textsuperscript{(73)}, \VUgras{cornetins} \textsuperscript{(74)}, \VUgras{clarinetes}
\textsuperscript{(75)}, violinos com os seus \VUgras{arcos} \textsuperscript{(76)} e até
um \VUgras{bombo} \textsuperscript{(77)}. Um pouco mais adiante, no balcão da
papelaria, um \VUgras{estudante} \textsuperscript{(78)} regateava com o
\VUgras{caixeiro} \textsuperscript{(79)} uma pasta de documentos. Perto deles
distinguia um belo \VUgras{abecedário} \textsuperscript{(80)}.

%%K t16-14-2 p
No meio dos armazéns erguia-se uma alta \VUgras{árvore de Natal}
\textsuperscript{(81)} carregada de velas, bugigangas e brinquedos de todas as
espécies: \VUgras{cavalos de madeira} \textsuperscript{(82)}, \VUgras{bonecas}
\textsuperscript{(83)}, etc.

%%K t16-14-3 p
A \VUgras{perfumaria} \textsuperscript{(84)}, com os seus \VUgras{dentífricos} e os
seus diversos \VUgras{perfumes}, exalava um cheiro muito agradável que
subia até nós.

%%K t16-tit-5 sub
\VUsaut{3.4mm}
\VUpk{10.2pt}{40}{Compramos alguma coisa.}
\VUsaut{2.8mm}

%%K t16-15-1 p
15. --- Como o meu avô começava a achar o tempo longo, voltámos a descer
pela grande escada. Deteve-se um momento diante de um \VUgras{quadro de
astronomia} \textsuperscript{(85)} que representava, segundo me disse,
\VUgras{cometas} \textsuperscript{(a)}, \VUgras{eclipses da lua e do sol}
\textsuperscript{(b)}, uma rosa dos ventos com os \VUgras{quatro pontos cardeais}
\textsuperscript{(c)} \VUgras{(norte, sul, este e oeste)}, um mapa dos dois
\VUgras{hemisférios} \textsuperscript{(d)}, os \VUgras{planetas} \textsuperscript{(e)} e o
\VUgras{sistema solar} \textsuperscript{(f)}. Não percebi nada daquilo, e
interessaram-me muito mais a libré galoada de um \VUgras{entregador}
\textsuperscript{(86)} que passava, e os bonitos \VUgras{leques}
\textsuperscript{(87)} que tinha ao lado, junto aos \VUgras{porta-moedas}
\textsuperscript{(88)} e às \VUgras{carteiras} \textsuperscript{(89)}.

%%K t16-16-1 p
16. --- Admirei também no balcão umas \VUgras{plumas} \textsuperscript{(90)} e umas
\VUgras{fivelas de cinto} \textsuperscript{(91)}, e as bonitas \VUgras{flores
artificiais} \textsuperscript{(94)} graciosamente dispostas em cestos. As
\VUgras{violetas}, os \VUgras{goivos}, os \VUgras{jacintos}
\textsuperscript{(95)}, os \VUgras{gerânios} \textsuperscript{(96)}, os \VUgras{lírios}
\textsuperscript{(97)} e as \VUgras{chagas} \textsuperscript{(98)} estavam muito bem
imitados.

%%K t16-tit-6 sub
\VUsaut{4.4mm}
\VUpk{10.2pt}{40}{O presente do avô.}
\VUsaut{2.8mm}

%%K t16-17-1 p
17. --- Pedi ao meu avô que me comprasse uma das \VUgras{escovas de dentes}
\textsuperscript{(92)} de uma caixa que havia junto aos \VUgras{ganchos de cabelo}
\textsuperscript{(93)}. Aceitou, e fui eu mesmo pagar a minha compra à caixa, ao
lado da qual estavam a fazer um grande \VUgras{embrulho} \textsuperscript{(100)}
para um \VUgras{cliente} \textsuperscript{(99)}.

%%K t16-18-1 p
18. --- De repente houve um ligeiro alvoroço entre a gente que estava junto
aos \VUgras{bronzes} \textsuperscript{(101)} artísticos. Um \VUgras{ladrão}
\textsuperscript{(102)} acabava de ser surpreendido em flagrante delito por um
\VUgras{inspetor dos armazéns} \textsuperscript{(103)} quando tentava roubar
alguém. Mandou-se buscar um polícia para o \VUgras{prender}, e no momento
em que saíamos, levavam o culpado para a esquadra.
\VUsaut{6.0mm}
\VUfilet{22mm}
